\chapter{Implementation and Testing}

This chapter provides a comprehensive overview of the implementation and testing approaches employed in developing the Urdu AI Video Creator platform. The system integrates modern authentication, AI-powered script generation, and content creation workflows to deliver a seamless user experience for creating multilingual video content.

% =====================================================
% GENERAL DESCRIPTION
% =====================================================

\section{General Description}

The Urdu AI Video Creator is a Next.js-based web application designed to enable content creators to generate AI-powered videos with support for Urdu and English languages. The platform features enterprise-grade authentication, intelligent script generation, voice cloning capabilities, and lip-sync technology.

\subsection{System Context}

The system addresses the growing demand for multilingual content creation tools, particularly for Urdu language content. Traditional video creation workflows are time-consuming and require significant technical expertise. This platform automates script generation, voice synthesis, and video production through AI integration.

\subsection{Technology Stack}

\begin{itemize}
    \item \textbf{Frontend}: Next.js 14 with TypeScript, Tailwind CSS, Radix UI components
    \item \textbf{Backend}: Node.js with Express.js, TypeScript
    \item \textbf{Authentication}: Clerk for enterprise-grade authentication and session management
    \item \textbf{Database}: Supabase (PostgreSQL) with Row Level Security
    \item \textbf{AI Services}: OpenRouter API for LLM integration (GPT-4), OpenAI Whisper for transcription
    \item \textbf{Text-to-Speech}: Coqui XTTS v2 for voice cloning and synthesis (English: pre-trained, Urdu: custom fine-tuned)
    \item \textbf{Additional Libraries}: Framer Motion for animations, PyTorch for model inference
\end{itemize}

\subsection{Current Implementation Status}

The current iteration includes:
\begin{itemize}
    \item Complete authentication system with user management
    \item Multi-stage content creation workflow
    \item AI-powered script generation and refinement
    \item Iterative feedback-based script improvement
    \item Draft management and version control
    \item Protected dashboard with user-specific content access
\end{itemize}

% =====================================================
% ALGORITHM DESIGN
% =====================================================

\section{Algorithm Design}

This section details the core algorithms and workflows implemented in the system.

\subsection{User Authentication and Synchronization}

The authentication flow ensures secure user access and data synchronization between Clerk and Supabase.

\begin{algorithm}[H]
\caption{User Authentication and Synchronization}
\begin{algorithmic}[1]
\State \textbf{Input:} User credentials (email, password) or OAuth provider
\State \textbf{Output:} Authenticated session with synchronized user data

\Procedure{AuthenticateUser}{credentials}
    \State user $\gets$ Clerk.authenticate(credentials)
    \If{authentication successful}
        \State token $\gets$ generateJWT(user.clerkId)
        \State session $\gets$ createSession(token)
        \State \Call{SyncUserToSupabase}{user.clerkId, user.email, user.firstName, user.lastName}
        \State \textbf{return} session
    \Else
        \State \textbf{return} error("Authentication failed")
    \EndIf
\EndProcedure

\Procedure{SyncUserToSupabase}{clerkId, email, firstName, lastName}
    \State token $\gets$ getCurrentSessionToken()
    \State headers $\gets$ \{"Authorization": "Bearer " + token\}
    \State body $\gets$ \{clerkId, email, firstName, lastName\}
    \State response $\gets$ POST("/api/users/sync", headers, body)
    \If{response.success}
        \State Supabase.upsert("users", \{clerk\_id: clerkId, email, first\_name: firstName, last\_name: lastName\})
    \EndIf
\EndProcedure

\Procedure{VerifyRequest}{request}
    \State token $\gets$ extractToken(request.headers["Authorization"])
    \State session $\gets$ Clerk.verifyToken(token)
    \If{session is valid}
        \State request.auth $\gets$ \{userId: session.sub\}
        \State \textbf{continue} to protected route
    \Else
        \State \textbf{return} 401 Unauthorized
    \EndIf
\EndProcedure
\end{algorithmic}
\end{algorithm}

\subsection{AI Script Generation Workflow}

The script generation algorithm leverages large language models to create contextually appropriate content based on user specifications.

\begin{algorithm}[H]
\caption{AI-Powered Script Generation}
\begin{algorithmic}[1]
\State \textbf{Input:} ScriptParameters (title, topic, language, scriptType, tone, targetAudience, duration, pacing, etc.)
\State \textbf{Output:} Generated script with metadata

\Procedure{GenerateScript}{params}
    \State userId $\gets$ getCurrentUserId()
    \State \Call{ValidateParameters}{params}
    
    \State prompt $\gets$ \Call{BuildGenerationPrompt}{params}
    \Comment{Construct structured prompt}
    
    \State messages $\gets$ [\{"role": "user", "content": prompt\}]
    \State generatedText $\gets$ \Call{CallLLM}{messages, maxTokens=4000}
    
    \State wordCount $\gets$ \Call{CountWords}{generatedText}
    \State estimatedDuration $\gets$ \Call{CalculateDuration}{wordCount, params.pacing}
    
    \State scriptId $\gets$ generateUUID()
    \State metadata $\gets$ \{wordCount, estimatedDuration, scriptType: params.scriptType, tone: params.tone\}
    
    \State script $\gets$ \{
        \State \quad scriptId, userId, title: params.title, language: params.language,
        \State \quad content: generatedText, method: "generated", metadata, status: "draft"
    \State \}
    
    \State \Call{SaveToDatabase}{script}
    
    \State \textbf{return} \{scriptId, content: generatedText, metadata\}
\EndProcedure

\Procedure{BuildGenerationPrompt}{params}
    \State template $\gets$ "You are a professional content writer..."
    \State template += "Language: " + params.language
    \State template += "Topic: " + params.topic
    \State template += "Script Type: " + params.scriptType
    \State template += "Tone: " + params.tone
    \State template += "Target Audience: " + params.targetAudience
    \State template += "Duration: " + params.duration + " seconds"
    \State template += "Pacing: " + params.pacing
    
    \If{params.includeHook}
        \State template += "Include an attention-grabbing hook in the first 5 seconds"
    \EndIf
    
    \If{params.includeCTA}
        \State template += "Include a clear call-to-action at the end"
    \EndIf
    
    \If{params.keyPoints}
        \State template += "Key points to cover: " + params.keyPoints
    \EndIf
    
    \State \textbf{return} template
\EndProcedure

\Procedure{CalculateDuration}{wordCount, pacing}
    \If{pacing == "slow"}
        \State wordsPerMinute $\gets$ 120
    \ElsIf{pacing == "medium"}
        \State wordsPerMinute $\gets$ 140
    \Else \Comment{fast}
        \State wordsPerMinute $\gets$ 160
    \EndIf
    \State duration $\gets$ wordCount / wordsPerMinute
    \State \textbf{return} duration \Comment{in minutes}
\EndProcedure
\end{algorithmic}
\end{algorithm}

\subsection{Iterative Script Refinement with Feedback}

This algorithm enables users to refine scripts through iterative feedback, maintaining version history.

\begin{algorithm}[H]
\caption{Feedback-Based Script Refinement}
\begin{algorithmic}[1]
\State \textbf{Input:} scriptId, currentScript, feedback, language, duration, pacing
\State \textbf{Output:} Refined script with incremented version number

\Procedure{RefineWithFeedback}{scriptId, currentScript, feedback, language, duration, pacing}
    \State userId $\gets$ getCurrentUserId()
    \State existingScript $\gets$ \Call{FetchScript}{scriptId, userId}
    
    \If{existingScript is null}
        \State \textbf{return} error("Script not found")
    \EndIf
    
    \State currentVersion $\gets$ length(existingScript.versions)
    
    \State prompt $\gets$ \Call{BuildFeedbackPrompt}{language, duration, pacing, feedback, currentVersion}
    \State prompt += "\textbackslash n\textbackslash nCurrent script:\textbackslash n" + currentScript
    
    \State messages $\gets$ [\{"role": "user", "content": prompt\}]
    \State refinedText $\gets$ \Call{CallLLM}{messages, maxTokens=4000}
    
    \State metadata $\gets$ \Call{CalculateMetadata}{refinedText, pacing}
    
    \State newVersion $\gets$ \{
        \State \quad versionNumber: currentVersion + 1,
        \State \quad content: refinedText,
        \State \quad feedback: feedback,
        \State \quad createdAt: getCurrentTimestamp()
    \State \}
    
    \State existingScript.versions.append(newVersion)
    \State existingScript.content $\gets$ refinedText
    \State existingScript.metadata $\gets$ metadata
    
    \State \Call{UpdateDatabase}{scriptId, existingScript}
    
    \State \textbf{return} \{scriptId, content: refinedText, metadata, version: currentVersion + 1\}
\EndProcedure

\Procedure{BuildFeedbackPrompt}{language, duration, pacing, feedback, version}
    \State prompt $\gets$ "You are refining a script based on user feedback."
    \State prompt += "This is version " + (version + 1) + " of the script."
    \State prompt += "Language: " + language
    \State prompt += "Target duration: " + duration + " seconds"
    \State prompt += "Pacing: " + pacing
    \State prompt += "\textbackslash n\textbackslash nUser feedback: " + feedback
    \State prompt += "\textbackslash n\textbackslash nApply the feedback to improve the script while maintaining its core message."
    \State \textbf{return} prompt
\EndProcedure
\end{algorithmic}
\end{algorithm}

\subsection{Script Refinement (Simple and Custom)}

\begin{algorithm}[H]
\caption{Script Refinement Process}
\begin{algorithmic}[1]
\State \textbf{Input:} originalScript, refinementType, customInstructions (optional), language, duration, pacing
\State \textbf{Output:} Refined script

\Procedure{RefineScript}{originalScript, refinementType, customInstructions, language, duration, pacing}
    \State userId $\gets$ getCurrentUserId()
    
    \If{refinementType == "simple"}
        \State prompt $\gets$ \Call{GetSimpleRefinementPrompt}{language, duration, pacing}
    \ElsIf{refinementType == "custom"}
        \If{customInstructions is empty}
            \State \textbf{return} error("Custom instructions required")
        \EndIf
        \State prompt $\gets$ \Call{GetCustomRefinementPrompt}{language, duration, pacing, customInstructions}
    \Else
        \State \textbf{return} error("Invalid refinement type")
    \EndIf
    
    \State prompt += "\textbackslash n\textbackslash nOriginal script:\textbackslash n" + originalScript
    
    \State messages $\gets$ [\{"role": "user", "content": prompt\}]
    \State refinedText $\gets$ \Call{CallLLM}{messages, maxTokens=4000}
    
    \State metadata $\gets$ \Call{CalculateMetadata}{refinedText, pacing}
    \State scriptId $\gets$ generateUUID()
    
    \State script $\gets$ \{scriptId, userId, content: refinedText, method: "refinement", metadata\}
    \State \Call{SaveToDatabase}{script}
    
    \State \textbf{return} \{scriptId, content: refinedText, metadata\}
\EndProcedure
\end{algorithmic}
\end{algorithm}

\subsection{Voice Cloning with XTTS v2}

The text-to-speech synthesis uses Coqui XTTS v2, a state-of-the-art voice cloning model supporting both zero-shot and few-shot approaches.

\begin{algorithm}[H]
\caption{Voice Cloning and Speech Synthesis}
\begin{algorithmic}[1]
\State \textbf{Input:} script text, language, speaker\_wav (reference audio for cloning)
\State \textbf{Output:} Synthesized speech audio file

\Procedure{SynthesizeSpeech}{text, language, speaker\_wav}
    \State device $\gets$ "cuda" if GPU available else "cpu"
    
    \If{language == "English"}
        \State modelName $\gets$ "tts\_models/multilingual/multi-dataset/xtts\_v2"
        \Comment{Pre-trained model}
    \ElsIf{language == "Urdu"}
        \State modelName $\gets$ "custom\_models/urdu/xtts\_v2\_finetuned"
        \Comment{Custom fine-tuned model}
    \Else
        \State \textbf{return} error("Unsupported language")
    \EndIf
    
    \State tts $\gets$ \Call{InitializeTTS}{modelName, device}
    
    \State \Call{ValidateReferenceAudio}{speaker\_wav}
    
    \State audioWav $\gets$ tts.tts(
        \State \quad text: text,
        \State \quad speaker\_wav: speaker\_wav,
        \State \quad language: language
    \State )
    
    \State outputPath $\gets$ generateOutputPath()
    \State tts.tts\_to\_file(
        \State \quad text: text,
        \State \quad speaker\_wav: speaker\_wav,
        \State \quad language: language,
        \State \quad file\_path: outputPath
    \State )
    
    \State \textbf{return} outputPath
\EndProcedure

\Procedure{ValidateReferenceAudio}{speaker\_wav}
    \If{not fileExists(speaker\_wav)}
        \State \textbf{return} error("Reference audio file not found")
    \EndIf
    \State duration $\gets$ getAudioDuration(speaker\_wav)
    \If{duration < 6 or duration > 30}
        \State \textbf{return} error("Reference audio should be 6-30 seconds")
    \EndIf
\EndProcedure
\end{algorithmic}
\end{algorithm}

\subsection{Custom XTTS v2 Training for Urdu Language}

Since no existing TTS model supports both Urdu language and voice cloning, a custom fine-tuning pipeline was developed to train XTTS v2 on Urdu language data.

\subsubsection{Dataset Collection and Preprocessing}

Multiple open-source Urdu speech datasets were collected and preprocessed:

\begin{itemize}
    \item \textbf{NLP-Story2Audio Dataset}: Urdu audio transcripts from GitHub (m01omer/NLP-Story2Audio)
    \item \textbf{Common Voice Urdu}: Preprocessed version by UmarRamzan (common-voice-urdu-processed)
    \item \textbf{Mendeley Urdu Speech Dataset}: Academic dataset (jcpfjnk5c2/3)
    \item \textbf{Jalandhary STT-TTS Dataset}: Kaggle dataset (mirfan899/jalandhary-stt-tts)
\end{itemize}

\begin{algorithm}[H]
\caption{Urdu Dataset Preprocessing Pipeline}
\begin{algorithmic}[1]
\State \textbf{Input:} Raw audio files and transcriptions from multiple sources
\State \textbf{Output:} Preprocessed dataset ready for training

\Procedure{PreprocessUrduDataset}{rawDatasets}
    \State combinedData $\gets$ []
    
    \For{each dataset in rawDatasets}
        \State audioFiles $\gets$ dataset.getAudioFiles()
        \State transcripts $\gets$ dataset.getTranscripts()
        
        \For{each (audio, transcript) pair}
            \State \Call{ValidateAudioQuality}{audio}
            \State \Call{NormalizeTranscript}{transcript}
            \State \Call{ResampleAudio}{audio, targetRate=22050}
            \State \Call{NormalizeAudioLevel}{audio}
            \State combinedData.append(\{audio, transcript\})
        \EndFor
    \EndFor
    
    \State trainSet, valSet $\gets$ \Call{SplitDataset}{combinedData, ratio=0.9}
    \State \textbf{return} trainSet, valSet
\EndProcedure

\Procedure{ValidateAudioQuality}{audio}
    \State sampleRate $\gets$ getAudioSampleRate(audio)
    \State duration $\gets$ getAudioDuration(audio)
    \State snr $\gets$ calculateSNR(audio)
    
    \If{duration < 1 or duration > 20}
        \State \textbf{skip} this sample
    \EndIf
    \If{snr < 20}
        \Comment{Signal-to-noise ratio threshold}
        \State \textbf{skip} this sample
    \EndIf
\EndProcedure

\Procedure{NormalizeTranscript}{transcript}
    \State transcript $\gets$ removeExtraWhitespace(transcript)
    \State transcript $\gets$ normalizeUrduCharacters(transcript)
    \State transcript $\gets$ removeDiacritics(transcript)
    \Comment{Optional}
    \State \textbf{return} transcript
\EndProcedure
\end{algorithmic}
\end{algorithm}

\subsubsection{Vocabulary Extension with BPE}

To handle Urdu script properly, the vocabulary was extended using Byte Pair Encoding (BPE).

\begin{algorithm}[H]
\caption{Vocabulary Extension for Urdu Script}
\begin{algorithmic}[1]
\State \textbf{Input:} Urdu text corpus from preprocessed transcripts
\State \textbf{Output:} Extended vocabulary tokenizer

\Procedure{ExtendVocabulary}{urduCorpus, baseTokenizer}
    \State allText $\gets$ concatenate all Urdu transcripts
    
    \State urduCharacters $\gets$ extractUniqueCharacters(allText)
    \State urduSubwords $\gets$ \Call{TrainBPE}{allText, vocabSize=5000}
    
    \State newVocab $\gets$ baseTokenizer.vocab
    \For{each token in urduSubwords}
        \If{token not in newVocab}
            \State newVocab.add(token)
        \EndIf
    \EndFor
    
    \State extendedTokenizer $\gets$ updateTokenizer(baseTokenizer, newVocab)
    \State \Call{SaveTokenizer}{extendedTokenizer, "urdu\_xtts\_tokenizer.json"}
    
    \State \textbf{return} extendedTokenizer
\EndProcedure

\Procedure{TrainBPE}{corpus, vocabSize}
    \State tokenizer $\gets$ ByteLevelBPETokenizer()
    \State tokenizer.train(
        \State \quad files: [corpus],
        \State \quad vocab\_size: vocabSize,
        \State \quad min\_frequency: 2,
        \State \quad special\_tokens: ["<pad>", "<sos>", "<eos>", "<unk>"]
    \State )
    \State \textbf{return} tokenizer.get\_vocab()
\EndProcedure
\end{algorithmic}
\end{algorithm}

\subsubsection{XTTS v2 Fine-tuning Pipeline}

The fine-tuning process adapts the pre-trained XTTS v2 model to Urdu language while preserving voice cloning capabilities.

\begin{algorithm}[H]
\caption{XTTS v2 Fine-tuning for Urdu}
\begin{algorithmic}[1]
\State \textbf{Input:} Pre-trained XTTS v2 model, preprocessed Urdu dataset, extended tokenizer
\State \textbf{Output:} Fine-tuned Urdu XTTS v2 model

\Procedure{FineTuneXTTS}{baseModel, urduDataset, tokenizer}
    \State config $\gets$ loadXTTSConfig("xtts\_v2\_config.json")
    \State config.vocab\_file $\gets$ "urdu\_xtts\_tokenizer.json"
    \State config.language $\gets$ "ur"
    
    \State model $\gets$ \Call{LoadPretrainedModel}{baseModel}
    \State model.tokenizer $\gets$ tokenizer
    
    \State \Comment{Freeze encoder layers, train only GPT and decoder}
    \For{layer in model.speaker\_encoder}
        \State layer.requires\_grad $\gets$ False
    \EndFor
    
    \State optimizer $\gets$ AdamW(
        \State \quad params: model.trainable\_parameters(),
        \State \quad lr: 5e-6,
        \State \quad weight\_decay: 1e-5
    \State )
    
    \State scheduler $\gets$ CosineAnnealingLR(optimizer, T\_max: epochs)
    
    \For{epoch in range(50)}
        \For{batch in urduDataset.train\_loader}
            \State audio, text, speaker\_embedding $\gets$ batch
            
            \State loss $\gets$ model.compute\_loss(audio, text, speaker\_embedding)
            \State loss.backward()
            \State optimizer.step()
            \State optimizer.zero\_grad()
            
            \If{step modulo 100 == 0}
                \State \Call{LogMetrics}{epoch, step, loss}
                \State \Call{ValidateModel}{model, urduDataset.val\_loader}
            \EndIf
        \EndFor
        
        \State scheduler.step()
        \State \Call{SaveCheckpoint}{model, epoch}
    \EndFor
    
    \State \textbf{return} model
\EndProcedure

\Procedure{ValidateModel}{model, valLoader}
    \State totalLoss $\gets$ 0
    \For{batch in valLoader}
        \State audio, text, speaker\_embedding $\gets$ batch
        \State loss $\gets$ model.compute\_loss(audio, text, speaker\_embedding)
        \State totalLoss += loss
    \EndFor
    \State avgLoss $\gets$ totalLoss / length(valLoader)
    \State \textbf{log} validation loss: avgLoss
\EndProcedure
\end{algorithmic}
\end{algorithm}

\subsubsection{Training Environment Challenges}

Due to Coqui TTS company shutdown in 2024, the library lacks support for latest Python versions and dependencies. The training was conducted on local systems with compatible environments:

\begin{itemize}
    \item \textbf{Python Version}: 3.9.x (compatibility requirement)
    \item \textbf{PyTorch}: 1.13.x with CUDA 11.7
    \item \textbf{Training Hardware}: NVIDIA GPU with minimum 16GB VRAM
    \item \textbf{Training Duration}: Approximately 48-72 hours on consumer-grade GPU
\end{itemize}

% =====================================================
% EXTERNAL APIs/SDKs
% =====================================================

\section{External APIs/SDKs}

The following table describes third-party APIs and SDKs integrated into the system:

\begin{table}[H]
\resizebox{\textwidth}{!}{%
\begin{tabular}{|p{4cm}|p{5cm}|p{4cm}|p{5cm}|}
\hline
\textbf{API and Version} & \textbf{Description} & \textbf{Purpose of Usage} & \textbf{API Endpoint/Function/Class Used} \\ \hline

Clerk (v6.35.6) & Enterprise authentication and user management & User signup, login, session management, JWT token generation & 
\texttt{ClerkProvider}, \texttt{useAuth()}, \texttt{clerkClient.verifyToken()}, \texttt{clerkMiddleware()} \\ \hline

Supabase (v2.86.2) & PostgreSQL database with real-time capabilities & Store user data, scripts, versions, and metadata with RLS & 
\texttt{createClient()}, \texttt{from('users')}, \texttt{from('scripts')}, \texttt{upsert()}, \texttt{select()} \\ \hline

OpenRouter API & LLM routing service for GPT-4 access & Generate and refine scripts using AI models & 
\texttt{POST https://openrouter.ai/api/v1/chat/completions} \\ \hline

OpenAI API (v5.20.3) & OpenAI services including Whisper & Audio transcription and language detection & 
\texttt{openai.audio.transcriptions.create()}, \texttt{whisper-1 model} \\ \hline

Coqui XTTS v2 & Advanced multilingual TTS with voice cloning & Zero-shot and few-shot voice cloning for English; custom fine-tuned model for Urdu language support & 
\texttt{TTS.api.TTS()}, \texttt{tts.tts\_to\_file()}, \texttt{xtts\_v2 model} \\ \hline

Next.js (v14.2.32) & React framework for full-stack applications & Server-side rendering, API routes, routing, and middleware & 
\texttt{next/navigation}, \texttt{middleware.ts}, API Routes \\ \hline

Radix UI (Multiple) & Accessible UI components library & Build accessible forms, dialogs, dropdowns, and interactive elements & 
\texttt{@radix-ui/react-dialog}, \texttt{@radix-ui/react-select}, \texttt{@radix-ui/react-checkbox} \\ \hline

Framer Motion (latest) & Animation library for React & Smooth transitions and animations in UI & 
\texttt{motion.div}, \texttt{AnimatePresence}, \texttt{variants} \\ \hline

\end{tabular}%
}
\caption{External APIs and SDKs used in the project}
\end{table}

% =====================================================
% TESTING DETAILS
% =====================================================

\section{Testing Details}

Testing ensures that the system functions correctly, securely, and efficiently. This section covers the testing methodologies employed during development.

\subsection{Unit Testing}

Unit tests validate individual components and functions in isolation. Each test targets a specific function to ensure correctness.

\subsubsection{Authentication Middleware Test}

\textbf{Test Case:} Verify that the authentication middleware correctly validates JWT tokens and rejects unauthorized requests.

\begin{verbatim}
// Test File: server/src/middleware/__tests__/auth.test.ts

import { requireAuth } from '../auth';
import { Request, Response, NextFunction } from 'express';
import { clerkClient } from '@clerk/clerk-sdk-node';

jest.mock('@clerk/clerk-sdk-node');

describe('Authentication Middleware', () => {
  let mockRequest: Partial<Request>;
  let mockResponse: Partial<Response>;
  let nextFunction: NextFunction;

  beforeEach(() => {
    mockRequest = {
      headers: {},
    };
    mockResponse = {
      status: jest.fn().mockReturnThis(),
      json: jest.fn(),
    };
    nextFunction = jest.fn();
  });

  test('should reject requests without Authorization header', async () => {
    await requireAuth(mockRequest as Request, mockResponse as Response, 
                     nextFunction);
    
    expect(mockResponse.status).toHaveBeenCalledWith(401);
    expect(mockResponse.json).toHaveBeenCalledWith({
      error: 'Missing or invalid Authorization header'
    });
    expect(nextFunction).not.toHaveBeenCalled();
  });

  test('should reject requests with invalid Bearer token', async () => {
    mockRequest.headers = { authorization: 'Bearer invalid_token' };
    
    (clerkClient.verifyToken as jest.Mock).mockRejectedValue(
      new Error('Invalid token')
    );
    
    await requireAuth(mockRequest as Request, mockResponse as Response, 
                     nextFunction);
    
    expect(mockResponse.status).toHaveBeenCalledWith(401);
    expect(mockResponse.json).toHaveBeenCalledWith({
      error: 'Authentication failed'
    });
  });

  test('should attach userId to request for valid tokens', async () => {
    const validToken = 'valid_jwt_token';
    mockRequest.headers = { authorization: `Bearer ${validToken}` };
    
    (clerkClient.verifyToken as jest.Mock).mockResolvedValue({
      sub: 'user_123456'
    });
    
    await requireAuth(mockRequest as Request, mockResponse as Response, 
                     nextFunction);
    
    expect(mockRequest.auth).toEqual({
      userId: 'user_123456',
      getToken: expect.any(Function)
    });
    expect(nextFunction).toHaveBeenCalled();
  });
});
\end{verbatim}

\textbf{Expected Results:}
\begin{itemize}
    \item Requests without authorization headers return 401 status
    \item Invalid tokens are rejected with appropriate error messages
    \item Valid tokens attach user ID to request and proceed to next middleware
\end{itemize}

\subsubsection{Script Generation Metadata Calculation Test}

\textbf{Test Case:} Validate that word count and duration calculations are accurate for different pacing settings.

\begin{verbatim}
// Test File: server/src/controllers/__tests__/contentController.test.ts

import { calculateMetadata } from '../contentController';

describe('Script Metadata Calculation', () => {
  
  test('should calculate correct word count', () => {
    const script = 'This is a sample script with ten words here.';
    const { wordCount } = calculateMetadata(script, 'medium');
    
    expect(wordCount).toBe(10);
  });

  test('should calculate duration for slow pacing (120 WPM)', () => {
    const script = 'word '.repeat(240); // 240 words
    const { estimatedDuration } = calculateMetadata(script, 'slow');
    
    expect(estimatedDuration).toBe(2); // 240 / 120 = 2 minutes
  });

  test('should calculate duration for medium pacing (140 WPM)', () => {
    const script = 'word '.repeat(280); // 280 words
    const { estimatedDuration } = calculateMetadata(script, 'medium');
    
    expect(estimatedDuration).toBe(2); // 280 / 140 = 2 minutes
  });

  test('should calculate duration for fast pacing (160 WPM)', () => {
    const script = 'word '.repeat(320); // 320 words
    const { estimatedDuration } = calculateMetadata(script, 'fast');
    
    expect(estimatedDuration).toBe(2); // 320 / 160 = 2 minutes
  });

  test('should handle empty scripts', () => {
    const script = '';
    const { wordCount, estimatedDuration } = calculateMetadata(script, 'medium');
    
    expect(wordCount).toBe(0);
    expect(estimatedDuration).toBe(0);
  });
});
\end{verbatim}

\textbf{Expected Results:}
\begin{itemize}
    \item Word count accurately reflects number of words in script
    \item Duration calculations adjust correctly based on pacing (slow: 120 WPM, medium: 140 WPM, fast: 160 WPM)
    \item Edge cases like empty scripts are handled without errors
\end{itemize}

\subsection{Integration Testing}

Integration tests validate the interaction between multiple components and external services.

\subsubsection{User Synchronization Flow Test}

\textbf{Test Case:} Verify complete user authentication and Supabase synchronization flow.

\begin{verbatim}
// Integration Test: User Authentication and Sync

describe('User Authentication and Synchronization', () => {
  test('should authenticate user and sync to Supabase', async () => {
    // Step 1: Authenticate with Clerk
    const signUpResponse = await clerk.signUp({
      emailAddress: 'testuser@example.com',
      password: 'SecurePass123!',
      firstName: 'Test',
      lastName: 'User'
    });
    
    expect(signUpResponse.status).toBe('complete');
    const userId = signUpResponse.createdUserId;
    
    // Step 2: Get authentication token
    const session = await clerk.sessions.getSessionList()[0];
    const token = await session.getToken();
    
    // Step 3: Call sync endpoint
    const syncResponse = await fetch('http://localhost:5000/api/users/sync', {
      method: 'POST',
      headers: {
        'Authorization': `Bearer ${token}`,
        'Content-Type': 'application/json'
      },
      body: JSON.stringify({
        email: 'testuser@example.com',
        firstName: 'Test',
        lastName: 'User'
      })
    });
    
    expect(syncResponse.status).toBe(200);
    const syncData = await syncResponse.json();
    expect(syncData.success).toBe(true);
    
    // Step 4: Verify in Supabase
    const { data: user } = await supabase
      .from('users')
      .select('*')
      .eq('clerk_id', userId)
      .single();
    
    expect(user).not.toBeNull();
    expect(user.email).toBe('testuser@example.com');
    expect(user.first_name).toBe('Test');
    expect(user.last_name).toBe('User');
  });
});
\end{verbatim}

\subsubsection{Script Generation End-to-End Test}

\textbf{Test Case:} Test complete script generation workflow from API request to database storage.

\begin{verbatim}
// Integration Test: Script Generation Workflow

describe('Script Generation Workflow', () => {
  let authToken: string;
  let userId: string;

  beforeAll(async () => {
    // Setup authenticated user
    const auth = await setupTestUser();
    authToken = auth.token;
    userId = auth.userId;
  });

  test('should generate script and save to database', async () => {
    const scriptParams = {
      title: 'Test Educational Video',
      language: 'English',
      topic: 'Introduction to AI',
      scriptType: 'Educational',
      tone: 'Professional',
      targetAudience: 'Beginners',
      duration: 120,
      pacing: 'medium',
      includeHook: true,
      includeCTA: true
    };
    
    // Call generation endpoint
    const response = await fetch('http://localhost:5000/api/content/generate-script', {
      method: 'POST',
      headers: {
        'Authorization': `Bearer ${authToken}`,
        'Content-Type': 'application/json'
      },
      body: JSON.stringify(scriptParams)
    });
    
    expect(response.status).toBe(200);
    const result = await response.json();
    
    // Verify response structure
    expect(result).toHaveProperty('scriptId');
    expect(result).toHaveProperty('content');
    expect(result).toHaveProperty('metadata');
    expect(result.metadata).toHaveProperty('wordCount');
    expect(result.metadata).toHaveProperty('estimatedDuration');
    
    // Verify database entry
    const { data: script } = await supabase
      .from('scripts')
      .select('*')
      .eq('script_id', result.scriptId)
      .single();
    
    expect(script).not.toBeNull();
    expect(script.user_id).toBe(userId);
    expect(script.title).toBe('Test Educational Video');
    expect(script.language).toBe('English');
    expect(script.method).toBe('generated');
    expect(script.status).toBe('draft');
  });
});
\end{verbatim}

\subsection{Manual Testing}

Manual testing was conducted to validate user interface functionality and user experience.

\subsubsection{Dashboard Navigation Test}

\textbf{Test Scenario:} User navigates through the dashboard and accesses different sections.

\textbf{Steps:}
\begin{enumerate}
    \item Log in with valid credentials
    \item Verify dashboard loads with user-specific data
    \item Click "Create Content" in sidebar
    \item Verify navigation to content creation page
    \item Verify active route highlighting in sidebar
    \item Test logout functionality
\end{enumerate}

\textbf{Expected Results:}
\begin{itemize}
    \item All navigation links function correctly
    \item Active routes are visually highlighted
    \item Protected routes redirect to login when not authenticated
    \item Session persists across page refreshes
\end{itemize}

\subsubsection{Script Generation User Flow Test}

\textbf{Test Scenario:} Complete script generation workflow from user perspective.

\textbf{Steps:}
\begin{enumerate}
    \item Navigate to "Create Content"
    \item Select language (English/Urdu)
    \item Choose "Generate from scratch"
    \item Fill in all required fields:
    \begin{itemize}
        \item Title: "Social Media Marketing Guide"
        \item Topic: "How to grow your Instagram following"
        \item Script type: Educational
        \item Tone: Conversational
        \item Target audience: Small business owners
        \item Duration: 2 minutes
        \item Pacing: Medium
    \end{itemize}
    \item Enable "Include hook" and "Include CTA"
    \item Click "Generate Script"
    \item Wait for loading indicator
    \item Verify generated script appears
    \item Review word count and estimated duration
    \item Test "Make it shorter" quick action
    \item Verify script updates
    \item Click "Save as Draft"
    \item Verify success message
\end{enumerate}

\textbf{Expected Results:}
\begin{itemize}
    \item All form fields validate correctly
    \item Loading states appear during generation
    \item Generated script matches language and parameters
    \item Quick actions modify script appropriately
    \item Draft saves successfully to database
    \item User can retrieve draft later from dashboard
\end{itemize}

\subsection{Security Testing}

\subsubsection{Authentication Security Test}

\textbf{Test Case:} Verify that protected routes reject unauthorized access.

\begin{verbatim}
// Security Test: Route Protection

describe('Route Security', () => {
  test('should block access to protected API without token', async () => {
    const response = await fetch('http://localhost:5000/api/content/generate-script', {
      method: 'POST',
      headers: { 'Content-Type': 'application/json' },
      body: JSON.stringify({ topic: 'test' })
    });
    
    expect(response.status).toBe(401);
  });

  test('should block access with expired token', async () => {
    const expiredToken = 'expired_jwt_token';
    
    const response = await fetch('http://localhost:5000/api/users/sync', {
      method: 'POST',
      headers: {
        'Authorization': `Bearer ${expiredToken}`,
        'Content-Type': 'application/json'
      }
    });
    
    expect(response.status).toBe(401);
  });

  test('should enforce user-specific data access (RLS)', async () => {
    // User A creates a script
    const userAToken = await getTokenForUser('userA');
    const scriptResponse = await createScript(userAToken, { title: 'A Script' });
    const scriptId = scriptResponse.scriptId;
    
    // User B attempts to access User A's script
    const userBToken = await getTokenForUser('userB');
    const accessResponse = await fetch(
      `http://localhost:5000/api/content/script/${scriptId}`,
      {
        headers: { 'Authorization': `Bearer ${userBToken}` }
      }
    );
    
    expect(accessResponse.status).toBe(404); // Not found (RLS prevents access)
  });
});
\end{verbatim}

\subsection{XTTS v2 Voice Cloning Testing}

\subsubsection{English Voice Cloning Test}

\textbf{Test Case:} Validate zero-shot voice cloning with pre-trained XTTS v2 model for English.

\begin{verbatim}
// Test File: TTS_Models/tests/test_english_xtts.py

import torch
from TTS.api import TTS
import pytest

class TestEnglishXTTS:
    @classmethod
    def setup_class(cls):
        cls.device = "cuda" if torch.cuda.is_available() else "cpu"
        cls.tts = TTS("tts_models/multilingual/multi-dataset/xtts_v2").to(cls.device)
    
    def test_zero_shot_voice_cloning(self):
        """Test zero-shot voice cloning with reference audio"""
        test_text = "This is a test of voice cloning technology."
        speaker_wav = "test_data/reference_speaker.wav"
        output_path = "test_output/cloned_voice.wav"
        
        # Perform TTS
        self.tts.tts_to_file(
            text=test_text,
            speaker_wav=speaker_wav,
            language="en",
            file_path=output_path
        )
        
        # Verify output file exists and has content
        assert os.path.exists(output_path)
        audio_duration = get_audio_duration(output_path)
        assert audio_duration > 0
        assert audio_duration < 10  # Should be reasonable length
    
    def test_voice_similarity(self):
        """Test voice similarity between reference and synthesized audio"""
        reference_wav = "test_data/reference_speaker.wav"
        synthesized_wav = "test_output/cloned_voice.wav"
        
        # Extract speaker embeddings
        ref_embedding = extract_speaker_embedding(reference_wav)
        synth_embedding = extract_speaker_embedding(synthesized_wav)
        
        # Calculate cosine similarity
        similarity = cosine_similarity(ref_embedding, synth_embedding)
        
        # Voice similarity should be high (> 0.75)
        assert similarity > 0.75
    
    def test_audio_quality(self):
        """Test synthesized audio quality metrics"""
        output_wav = "test_output/cloned_voice.wav"
        
        # Load audio
        audio, sr = librosa.load(output_wav, sr=22050)
        
        # Check signal-to-noise ratio
        snr = calculate_snr(audio)
        assert snr > 20  # Minimum acceptable SNR
        
        # Check for clipping
        max_amplitude = np.max(np.abs(audio))
        assert max_amplitude < 0.99  # No clipping
\end{verbatim}

\textbf{Expected Results:}
\begin{itemize}
    \item Synthesized audio file generated successfully
    \item Voice similarity score > 0.75 (high similarity to reference)
    \item Audio quality metrics within acceptable ranges (SNR > 20dB, no clipping)
\end{itemize}

\subsubsection{Urdu Fine-tuned Model Test}

\textbf{Test Case:} Validate custom fine-tuned XTTS v2 model for Urdu language with voice cloning.

\begin{verbatim}
// Test File: TTS_Models/tests/test_urdu_xtts.py

import torch
from TTS.api import TTS
import pytest

class TestUrduXTTS:
    @classmethod
    def setup_class(cls):
        cls.device = "cuda" if torch.cuda.is_available() else "cpu"
        # Load custom fine-tuned model
        cls.tts = TTS(model_path="custom_models/urdu/xtts_v2_finetuned",
                      config_path="custom_models/urdu/config.json").to(cls.device)
    
    def test_urdu_text_synthesis(self):
        """Test Urdu text-to-speech synthesis"""
        urdu_text = "یہ آواز کی نقل کی ٹیکنالوجی کا ایک ٹیسٹ ہے"
        speaker_wav = "test_data/urdu_reference_speaker.wav"
        output_path = "test_output/urdu_synthesized.wav"
        
        # Perform TTS
        self.tts.tts_to_file(
            text=urdu_text,
            speaker_wav=speaker_wav,
            language="ur",
            file_path=output_path
        )
        
        assert os.path.exists(output_path)
        audio_duration = get_audio_duration(output_path)
        assert audio_duration > 0
    
    def test_urdu_tokenization(self):
        """Test proper tokenization of Urdu script"""
        urdu_text = "سلام، یہ ایک ٹیسٹ ہے"
        
        # Tokenize using extended vocabulary
        tokens = self.tts.tokenizer.encode(urdu_text)
        
        # Verify no unknown tokens
        unk_token_id = self.tts.tokenizer.token_to_id("<unk>")
        assert unk_token_id not in tokens
        
        # Verify tokens can be decoded back
        decoded_text = self.tts.tokenizer.decode(tokens)
        assert decoded_text == urdu_text
    
    def test_urdu_voice_cloning_quality(self):
        """Test quality of Urdu voice cloning"""
        urdu_text = "یہ آواز کی نقل کا ایک نمونہ ہے"
        speaker_wav = "test_data/urdu_reference_speaker.wav"
        output_path = "test_output/urdu_cloned.wav"
        
        self.tts.tts_to_file(
            text=urdu_text,
            speaker_wav=speaker_wav,
            language="ur",
            file_path=output_path
        )
        
        # Verify speaker similarity
        ref_embedding = extract_speaker_embedding(speaker_wav)
        synth_embedding = extract_speaker_embedding(output_path)
        similarity = cosine_similarity(ref_embedding, synth_embedding)
        
        # Urdu model should maintain voice similarity > 0.70
        assert similarity > 0.70
    
    def test_dataset_preprocessing_validation(self):
        """Validate preprocessed Urdu dataset quality"""
        dataset_path = "TTS_Models/datasets/urdu_processed"
        
        # Load dataset metadata
        metadata = load_metadata(f"{dataset_path}/metadata.csv")
        
        # Check minimum samples
        assert len(metadata) > 1000
        
        # Verify audio files exist
        for row in metadata[:100]:  # Sample check
            audio_path = row['audio_file']
            assert os.path.exists(audio_path)
            
            # Check audio duration
            duration = get_audio_duration(audio_path)
            assert 1 <= duration <= 20
    
    def test_model_inference_speed(self):
        """Test inference speed for real-time requirements"""
        urdu_text = "یہ ایک مختصر جملہ ہے"
        speaker_wav = "test_data/urdu_reference_speaker.wav"
        
        import time
        start_time = time.time()
        
        wav = self.tts.tts(
            text=urdu_text,
            speaker_wav=speaker_wav,
            language="ur"
        )
        
        inference_time = time.time() - start_time
        audio_duration = len(wav) / 22050  # Assuming 22050 Hz sample rate
        
        # Real-time factor should be < 1.0 for acceptable performance
        rtf = inference_time / audio_duration
        assert rtf < 2.0  # Acceptable for non-real-time applications
\end{verbatim}

\textbf{Expected Results:}
\begin{itemize}
    \item Urdu text synthesized correctly without tokenization errors
    \item Voice cloning maintains similarity > 0.70 for Urdu speakers
    \item Dataset preprocessing validates minimum quality standards
    \item Inference time acceptable for production use (RTF < 2.0)
\end{itemize}

\subsection{Test Coverage Summary}

The testing strategy covers:

\begin{itemize}
    \item \textbf{Unit Tests}: Individual functions for authentication, metadata calculation, prompt generation, tokenization
    \item \textbf{Integration Tests}: End-to-end workflows including user sync and script generation
    \item \textbf{Manual Tests}: UI/UX validation, user flows, cross-browser compatibility
    \item \textbf{Security Tests}: Authentication enforcement, authorization checks, data isolation
    \item \textbf{TTS Model Tests}: Voice cloning quality, Urdu language synthesis, dataset validation, inference performance
\end{itemize}

All tests are designed to ensure system reliability, security, and correct functionality across different usage scenarios. Automated tests can be run using:

\begin{verbatim}
# Run all tests
npm test

# Run tests with coverage report
npm run test:coverage

# Run specific test file
npm test auth.test.ts

# Run XTTS model tests
cd TTS_Models
python -m pytest tests/ -v

# Run Urdu model specific tests
python -m pytest tests/test_urdu_xtts.py -v
\end{verbatim}

% =====================================================
% CONCLUSION
% =====================================================

\section{Conclusion}

The implementation follows industry best practices for authentication, database design, and API architecture. The modular structure allows for easy maintenance and future enhancements. Testing validates that core functionalities work correctly and securely.

A significant achievement of this iteration is the successful development of a custom XTTS v2 model for Urdu language support. By collecting and preprocessing multiple open-source Urdu speech datasets, extending the vocabulary using Byte Pair Encoding, and fine-tuning the pre-trained XTTS v2 model, the system now supports high-quality voice cloning in both English and Urdu languages. This addresses a critical gap in multilingual TTS systems, as no existing commercial solution provided both Urdu language support and voice cloning capabilities.

For English language synthesis, the system leverages the pre-trained XTTS v2 model's zero-shot and few-shot voice cloning capabilities, enabling users to clone voices with minimal reference audio. For Urdu, the custom fine-tuned model maintains comparable voice similarity scores (>0.70) while properly handling Urdu script and phonetics.

The challenges encountered during training—particularly regarding library compatibility due to Coqui TTS company shutdown in 2024—were successfully mitigated by establishing a stable local training environment with compatible Python and PyTorch versions.

The current iteration successfully delivers a functional content creation platform with AI-powered script generation and multilingual voice synthesis, setting a solid foundation for subsequent features including lip-sync video generation and complete end-to-end content production workflows.
